\documentclass[10pt,reqno]{amsart}

\usepackage[T1]{fontenc}
\usepackage[utf8]{inputenc}
\usepackage{newtxtext,newtxmath}
\usepackage{microtype}
\usepackage{mathtools}
\usepackage{booktabs}
\usepackage{enumitem}
\usepackage[nameinlink,noabbrev]{cleveref}
\usepackage[a4paper,margin=28mm,headheight=15pt]{geometry}

\setlength{\parskip}{2pt}
\setlist[enumerate]{leftmargin=7mm,itemsep=2pt,topsep=3pt}
\allowdisplaybreaks
\emergencystretch=3em

\newtheorem{theorem}{Theorem}[section]
\newtheorem{proposition}[theorem]{Proposition}
\newtheorem{lemma}[theorem]{Lemma}
\newtheorem{corollary}[theorem]{Corollary}
\theoremstyle{definition}
\newtheorem{definition}[theorem]{Definition}
\newtheorem{remark}[theorem]{Remark}

\newcommand{\C}{\mathbb C}
\newcommand{\D}{\mathbb D}
\newcommand{\cC}{\mathcal C}
\newcommand{\cR}{\mathcal R}
\newcommand{\hm}{\omega}

\title[Positive harmonic measure for a translated-cusp set]
{A Translated-Cusp Boundary Set\\of Positive Harmonic Measure}

\author{Argus Open Research Project}
\address{Argus Open Research Project}
\date{August 31, 2026}

\subjclass[2020]{Primary 30C85; Secondary 30C20, 31A15}
\keywords{harmonic measure, Jordan domain, boundary cusp, conformal mapping, open problem}

\begin{document}

\begin{abstract}
Problem~7.15 of the Canterbury problem collection asks for which
\(\alpha\in(1/2,1)\) a boundary set defined by local containment in
translates of
\[
 \cC_\alpha=\{x+iy:0<x<|y|^\alpha\}
\]
must have harmonic measure zero.  The printed replacement clause does not
separately specify how the distinguished vertex of the preceding angular
definition is transported.  We therefore state the conclusion relative to
an explicit source-vocabulary convention: the single comparison object may
be replaced, retained, weakened to boundary contact, or expanded at the same
parameter, but no new atomic predicate and no second-parameter copy may be
introduced.  Under every endpoint-consistent normal form in this convention,
we give a bounded Jordan domain for which the resulting set at
\(\alpha=2/3\) has positive harmonic measure from every interior pole.  The
construction is exact.  The note makes no claim of priority and records the
interpretive boundary explicitly.
\end{abstract}

\maketitle

\section{Statement and scope}

For \(0<\alpha<1\), put
\[
 \cC_\alpha=\{X+iY:0<X<|Y|^\alpha\}.
\]
This set has upper and lower connected components \(\cC_\alpha^+\) and
\(\cC_\alpha^-\), distinguished by the sign of \(Y\).

The source definition begins with angular boundary points.  A finite point
\(z_0\in\partial D\) is angular when some \(\varepsilon>0\) has the property
that every component \(U\) of \(D\cap B(z_0,\varepsilon)\) with
\(z_0\in\partial U\) is contained in an angle of aperture less than \(\pi\)
whose vertex is \(z_0\).  Problem~7.15 then replaces the angles by translates
of \(\cC_\alpha\).  The sentence identifies the new comparison set but does
not explicitly say what becomes of the vertex modifier.

We use the following finite interpretation class.  It is deliberately wider
than bare substitution, but it does not permit unrelated extra conditions.

\begin{definition}[source-vocabulary readings]\label{def:readings}
A reading in \(\cR_\alpha\) is obtained from the angular definition as follows.
\begin{enumerate}
\item Retain the tested point, the local radius, the local-component
      quantifier, the condition \(z_0\in\partial U\), and containment.
\item Replace the single angle comparator by a translate of
      \(\cC_\alpha\), \(\cC_\alpha^+\), or \(\cC_\alpha^-\).
\item Delete the angle vertex condition, retain it as equality with the cusp
      tip, or weaken it canonically to boundary contact.
\item Allow equivalent definitions at the same value of \(\alpha\), including
      any nonempty union of the canonical closure cells in
      \cref{eq:cells} below.
\item Allow the translating witness to be componentwise, common at one tested
      point, or common on the positive-measure witness set.
\end{enumerate}
No new atomic predicate on \(z_0-w\) is allowed, and the unique comparator
occurrence may not be copied, specialized at another parameter, and conjoined
back into the generic definition.
\end{definition}

Call a reading \emph{endpoint-consistent} if it does not contradict the
printed assertion that the zero-measure conclusion is false for a disk at
\(\alpha=1/2\).

\begin{theorem}\label{thm:main}
There is a bounded Jordan domain \(D\subset\C\) such that, for every
endpoint-consistent reading in \(\cR_{2/3}\), the corresponding boundary set
has strictly positive harmonic measure from every pole \(q\in D\).
In particular, under the source-vocabulary convention of
\cref{def:readings}, the universal zero-harmonic-measure property in
Problem~7.15 fails at \(\alpha=2/3\).
\end{theorem}

\begin{remark}\label{rem:scope}
The theorem is conditional on the stated interpretation class.  The source
does not explicitly name the contact point of the translated cusp.  A rule
that adds an unprinted proper-sublocus condition is outside
\(\cR_\alpha\) and can change the answer.  Thus \cref{thm:main} should not be
quoted without its source-vocabulary qualification.
\end{remark}

\section{The explicit Jordan domain}

Fix \(\alpha=2/3\), and write \(\cC=\cC_{2/3}\).  Define
\[
\begin{aligned}
 U_+&=\{x+iy:1<y<3,\ 0<x<y^{2/3}\},\\
 U_-&=\{x+iy:-3<y<-1,\ 0<x<(-y)^{2/3}\},\\
 R&=\{x+iy:1/4<x<1/2,\ -2<y<2\},\\
 D&=U_+\cup R\cup U_-.
\end{aligned}
\tag{2.1}\label{eq:domain}
\]
The point \(p=3/8\) belongs to \(D\).

\begin{proposition}\label{prop:jordan}
The set \(D\) is a bounded Jordan domain.  Its boundary has positive
logarithmic capacity.
\end{proposition}

\begin{proof}
The three sets in \cref{eq:domain} are open and connected, and
\[
 U_+\cap R=\{1<y<2,\ 1/4<x<1/2\},\qquad
 U_-\cap R=\{-2<y<-1,\ 1/4<x<1/2\}.
\]
Hence their union is a domain.  Its horizontal sections are
\[
 D_y=
 \begin{cases}
 (0,|y|^{2/3}),&1<|y|<3,\\
 (1/4,1/2),&|y|\leq1,\\
 \varnothing,&|y|\geq3.
 \end{cases}
\tag{2.2}\label{eq:sections}
\]
For \(1<|y|<2\), the rectangular section is already contained in the
cusp section.  Formula \cref{eq:sections} exhibits the boundary as a simple
closed chain consisting of four vertical lobe edges, two horizontal lobe
caps, two curved graphs \(x=|y|^{2/3}\), and the four sides joining the
central bridge to the lobes.  Nonconsecutive pieces are disjoint because the
curved pieces have \(x\geq1\), the bridge sides have \(x=1/4\) or \(1/2\),
and the lobe sides have \(x=0\), with disjoint ordinate ranges where needed.
Thus \(D\) is the bounded complementary component of a simple closed curve.

The segment \([i,3i]\) lies in \(\partial D\).  A line segment of length
\(L\) has logarithmic capacity \(L/4\); consequently
\(\operatorname{cap}(\partial D)\geq1/2\).
\end{proof}

\section{Four positive-measure cusp arcs}

For \(4/3<u<7/5\), set
\[
\begin{aligned}
 E_{v,+}&=\{iu^3\},&
 E_{c,+}&=\{u^2+iu^3\},\\
 E_{v,-}&=\{-iu^3\},&
 E_{c,-}&=\{u^2-iu^3\}.
\end{aligned}
\tag{3.1}\label{eq:arcs}
\]
These are open subarcs of \(\partial D\).  The subscripts \(v\) and \(c\)
refer to the vertical and curved boundary strata of \(\overline\cC\).

\begin{lemma}[uniform local equality]\label{lem:local}
For every \(z_0\) in any of the four arcs in \cref{eq:arcs},
\[
 D\cap B(z_0,1/8)=\cC\cap B(z_0,1/8).
\tag{3.2}\label{eq:local}
\]
\end{lemma}

\begin{proof}
The curved points lie on \(x=|y|^{2/3}\), since
\((u^2)^3=(u^3)^2\).  The vertical points lie on \(x=0\), with nonzero
imaginary part.  Moreover,
\[
 \frac{64}{27}<|\operatorname{Im}z_0|<\frac{343}{125}.
\]
If \(|z-z_0|<1/8\), then \(2<|\operatorname{Im}z|<3\), with the same sign
as \(\operatorname{Im}z_0\).  On these two strips the bridge is absent and
the appropriate lobe of \(D\) is exactly the corresponding component of
\(\cC\).  This proves \cref{eq:local}.
\end{proof}

\begin{proposition}\label{prop:hm}
Each of the four arcs in \cref{eq:arcs} has strictly positive harmonic
measure from every pole \(q\in D\).
\end{proposition}

\begin{proof}
Let \(\phi:D\to\D\) be a Riemann map.  By the Carath\'eodory boundary
theorem, \(\phi\) extends to a homeomorphism of the closures.  Every arc
\(E\) in \cref{eq:arcs} is nonempty and open in the Jordan curve
\(\partial D\), so \(J=\phi(E)\) is a nonempty open circular arc.  Conformal
invariance and the disk Poisson kernel give
\[
 \hm_D^q(E)=\frac1{2\pi}\int_J
 \frac{1-|\phi(q)|^2}{|\zeta-\phi(q)|^2}\,|d\zeta|>0.
\]
The integrand is positive and \(J\) has positive arclength.
\end{proof}

\section{Normal forms for the replacement clause}

For a translate \(w+\cC_\alpha\), write
\(z_0-w=X+iY\).  The closure of \(\cC_\alpha\) has the seven canonical
cells
\[
\begin{array}{c|c}
 O&X=0,\ Y=0\\
 V_+&X=0,\ Y>0\\
 V_-&X=0,\ Y<0\\
 \Gamma_+&X=Y^\alpha,\ Y>0\\
 \Gamma_-&X=(-Y)^\alpha,\ Y<0\\
 I_+&0<X<Y^\alpha,\ Y>0\\
 I_-&0<X<(-Y)^\alpha,\ Y<0.
\end{array}
\tag{4.1}\label{eq:cells}
\]
The full shape is compatible with all seven cells.  The upper branch is
compatible with \(O,V_+,\Gamma_+,I_+\), and the lower branch with
\(O,V_-,\Gamma_-,I_-\).

\begin{lemma}[finite normal-form count]\label{lem:count}
Definition~\ref{def:readings} has \(157\) nonempty compatible
shape-selector pairs.  With componentwise, pointwise-common, and
witness-set-common translate scope, it has \(471\) scope-labelled normal
forms.
\end{lemma}

\begin{proof}
For the full shape, every nonempty union of the seven cells is allowed,
giving \(2^7-1=127\) selectors.  Each branch has four compatible cells and
therefore \(2^4-1=15\) selectors.  Thus there are
\(127+15+15=157\) shape-selector pairs.  The three translate scopes multiply
this number by three.
\end{proof}

The endpoint assertion at \(\alpha=1/2\) restricts which selectors can
represent the printed definition.  Interior contact admits every point of a
disk after a sufficiently distant translate.  Curved contact admits a
nondegenerate disk arc by the standard tangent-parabola computation.  In
contrast, origin contact makes each connected local component lie in one
quadrant and hence makes the point angular; such points are harmonic-null.
Vertical-only contact on a disk occurs only at an extremal point and is also
harmonic-null.  Consequently an endpoint-consistent selector contains at
least one of
\[
 I_+,\ I_-,\ \Gamma_+,\ \Gamma_-.
\tag{4.2}\label{eq:survivors}
\]

\section{One domain handles every surviving reading}

The arcs in \cref{eq:arcs}, together with \(w=0\), handle the four
nonorigin boundary cells:
\[
 V_+,\quad V_-,\quad\Gamma_+,\quad\Gamma_-.
\]
By \cref{lem:local}, the same local radius \(1/8\) and the same translate
work at every point of each respective arc.

For the two interior cells, put
\[
 w_+=-1-100i,\qquad w_-=-1+100i.
\tag{5.1}\label{eq:far}
\]
For \(z=x+iy\in\overline D\), relative to either translate one has
\[
 1\leq X=x+1<4,\qquad |Y|\geq97>8=4^{3/2}.
\]
Hence \(X<|Y|^{2/3}\).  The sign of \(Y\) is positive relative to \(w_+\)
and negative relative to \(w_-\), so
\[
 \overline D\subset w_++I_+,
 \qquad
 \overline D\subset w_-+I_-.
\tag{5.2}\label{eq:interior}
\]
Thus a single translate handles every boundary point whenever the selector
contains the corresponding interior cell.

\begin{proof}[Proof of \cref{thm:main}]
Let an endpoint-consistent reading in \(\cR_{2/3}\) be given.  By
\cref{eq:survivors}, its selector contains an interior or curved cell.  If it
contains \(\Gamma_+\) or \(\Gamma_-\), choose the corresponding curved arc in
\cref{eq:arcs}; \cref{lem:local,prop:hm} provide local containment and
positive harmonic measure.  If it contains \(I_+\) or \(I_-\), use
\cref{eq:interior}; in this case all of \(\partial D\) is admitted.  The
argument uses a common translate, so it also satisfies every weaker
componentwise or pointwise scope.  Therefore every endpoint-consistent normal
form produces a boundary set of positive harmonic measure.
\end{proof}

\section{Why unrestricted enrichment is different}

The qualification in \cref{rem:scope} is substantive.  Consider
\[
 Q(X,Y):\quad 0<|Y|<1,\qquad X^2=|Y|.
\tag{6.1}\label{eq:q}
\]
The condition \(0<|Y|\) can be formed after selecting nonorigin cells, but
the cutoff \(|Y|<1\) is not supplied by the source.  The equation
\(X^2=|Y|\) is the curved-boundary equation only at \(\alpha=1/2\); at
\(\alpha=2/3\) the corresponding equation is \(X^3=|Y|^2\).  Conjoining both
requires copying the unique comparator occurrence and specializing the copy
at a second parameter.  Hence \(Q\) is excluded by
\cref{def:readings}.

This exclusion cannot be replaced by the vague assertion that \(Q\) is
unnatural.  If arbitrary pragmatic additions are allowed, \(Q\) is a genuine
endpoint-compatible rule that rejects the four frozen arcs, since their
heights satisfy
\[
 \frac{64}{27}<|Y|<\frac{343}{125}
\]
and therefore exceed one.  The theorem concerns the finite
occurrence-preserving, same-parameter source vocabulary, not unrestricted
enrichment.

\section{Verification and status}

Two independent exact checks were performed.  The geometric check verifies
the strict parameter range, connected overlaps, horizontal sections, Jordan
boundary, four arc types, local radius, curved-boundary identities, and the
distant-translate inequalities using integer and rational arithmetic.  The
source-form check regenerates all \(157\) compatible selectors and all
\(471\) scope-labelled normal forms, classifies the endpoint survivors, and
checks that every surviving cell has a witness in \cref{eq:arcs,eq:far}.

The mathematical conclusion and the source-vocabulary qualification passed
independent review.  No claim is made that the interpretation is forced by
an explicit contact convention absent from the printed sentence.  No claim
of novelty or priority is made.  The value of the note is the exact domain,
the robust finite interpretation audit, and the explicit statement of the
remaining semantic boundary.

\begin{thebibliography}{9}

\bibitem{ClunieHayman}
J.~G. Clunie and W.~K. Hayman, editors,
\emph{Proceedings of the Symposium on Complex Analysis, Canterbury 1973},
London Mathematical Society Lecture Note Series, vol.~12,
Cambridge University Press, 1974, pp.~155--180.

\bibitem{Pommerenke}
Ch. Pommerenke,
\emph{Boundary Behaviour of Conformal Maps},
Grundlehren der mathematischen Wissenschaften, vol.~299,
Springer-Verlag, Berlin, 1992.

\bibitem{Ransford}
T. Ransford,
\emph{Potential Theory in the Complex Plane},
London Mathematical Society Student Texts, vol.~28,
Cambridge University Press, 1995.

\end{thebibliography}

\end{document}
